\SourceLeaf{281}%
\(x\). Diese Koordinaten sind in Teilen des Radius zu verstehen. Man muss sie also noch mit \(A\) multiplizieren, wenn man sie als Längen angeben will, oder mit \(206265''\), wenn man sie im Bogenmaß angeben will.

Ein Element von \(M\) wird durch
\[
\BareRoot{(\cos y^2\mathrm{d}x^2+\mathrm{d}y^2)}
\]
oder durch \(\frac{\cos y}{\cos\psi}\mathord{.}\mathrm{d}x\) ausgedrückt, wenn man
\[
\frac{\mathrm{d}y}{\cos y\mathrm{d}x}=\operatorname{tang}\psi
\]
setzt. Dabei bezeichnet \(\psi\) also die Neigung des Elements gegenüber der Parallelen zu \(N\). Zur Veranschaulichung kann man sich vorstellen, dass die \(x\) von rechts nach links und die \(y\) von unten nach oben wachsen. Damit ist die Richtung positiver \(\psi\) von selbst bestimmt.

Das wie oben mit \(m\) bezeichnete Vergrößerungsverhältnis bei der Abbildung der Ellipsoidfläche auf die Kugelfläche lässt sich hier als Funktion von \(x\) und \(y\) auffassen. Die Länge des Elements von \(L\), das jenem Element von \(M\) entspricht, ist
\[
=\frac{A\cos y}{m\cos\psi}\mathord{.}\mathrm{d}x
\]
Setzt man zur Abkürzung
\[
\begin{aligned}
\log\operatorname{tang}(45^0+\tfrac{1}{2}y)&=u\\
\frac{\cos y}{m}&=n
\end{aligned}
\]
so ist auch \(n\) eine Funktion von \(x\) und \(y\) oder, was auf dasselbe hinausläuft, von \(x\) und \(u\). Damit hat man
\[
\operatorname{tang}\psi=\frac{\mathrm{d}u}{\mathrm{d}x}
\]
und das Element von \(L\)
\[
=\frac{An}{\cos\psi}\mathord{.}\mathrm{d}x
\]
\indent Die Beschaffenheit der Linie \(M\) wird also durch die Bedingung bestimmt, dass zwischen beliebigen vorgegebenen Grenzen das Integral \(\int\frac{n}{\cos\psi}\mathrm{d}x\) oder
\[
\int n\BareRoot{(1+\frac{\mathrm{d}u^2}{\mathrm{d}x^2})}\mathrm{d}x
\]
